%% it/sections/04_playbook.tex
%% ============================================================
%% 04_playbook.tex — 28 playbook (7 per fazione)
%% ============================================================

\chapter{I Playbook / The Playbooks}
\markboth{I Playbook}{The Playbooks}

\lettrine[lines=3]{O}{gni} personaggio di \textit{Vespri 1282} è definito
da un \textbf{playbook}: un archetipo storico che ne fissa il ruolo
narrativo, le azioni di punta, e la tensione centrale. I playbook non
sono classi rigide — sono punti di partenza per costruire una persona.

Ogni playbook prevede:
\begin{itemize}[nosep]
  \item \textbf{Azioni di punta} — due azioni con valore iniziale 2
  \item \textbf{Azioni secondarie} — quattro azioni con valore iniziale 1
  \item \textbf{Dono} — un'abilità speciale unica al playbook
  \item \textbf{Fardello} — una debolezza o complicazione strutturale
  \item \textbf{Legame iniziale} — un legame pre-scritto con un altro PG
    della stessa fazione
\end{itemize}

\secsep

%% ============================================================
%% FAZIONE I: CONGIURATI
%% ============================================================
\section{I Congiurati — Playbook}

\subsection{Il Messaggero}
{\cinzel\footnotesize\color{grigionote} Ruggero di Amato}\\
\textit{Porta parole che possono cambiare il mondo. O che possono
ucciderti.}

\textbf{Azioni di punta:} Manovrare, Ingannare\\
\textbf{Azioni secondarie:} Studiare, Agire, Persuadere, Resistere

\textbf{Dono — Volti multipli:} Hai almeno tre identità false pronte.
Quando dichiari di usarne una, il tuo tiro di Ingannare è sempre
almeno in posizione Controllata, a meno che non ti abbiano già visto
con quella faccia.

\textbf{Fardello — La catena:} Ogni messaggio che porti ti vincola
al suo mittente. Se vieni catturato, il GM può chiederti un tiro di
Resistere (Disperata) per non rivelare la fonte. Se fallisci, la rete
perde 1 punto.

\textbf{Legame iniziale:} Il Cospiratore sa il tuo vero nome. Sei
l'unico che lo sa.

\filetto

\subsection{Il Cospiratore}
{\cinzel\footnotesize\color{grigionote} Gualtiero de Caltagirone}\\
\textit{Conosci il piano. O almeno, la tua parte di esso.}

\textbf{Azioni di punta:} Tramare, Studiare\\
\textbf{Azioni secondarie:} Ingannare, Persuadere, Comandare, Resistere

\textbf{Dono — Compartimentazione:} Una volta per sessione, puoi
dichiarare che una cosa che stai facendo era «parte del piano»
fin dall'inizio, anche se non era stata esplicitata. Questo non
cambia i dadi — cambia la narrazione.

\textbf{Fardello — Il peso di sapere:} Sai cose che gli altri non sanno.
Questo ti isola. Quando un altro PG della tua fazione tira Studiare
per capire la situazione e fallisce, tu puoi scegliere di rivelargli
l'informazione — ma subisci 1 Stress per averlo fatto.

\textbf{Legame iniziale:} Il Messaggero porta le tue parole. Speri
che sia abbastanza prudente.

\filetto

\subsection{La Spia}
{\cinzel\footnotesize\color{grigionote} Margherita di Brindisi}\\
\textit{Lavori dentro il sistema angioino. Da troppo tempo.}

\textbf{Azioni di punta:} Ingannare, Studiare\\
\textbf{Azioni secondarie:} Manovrare, Agire, Persuadere, Resistere

\textbf{Dono — Dentro il palazzo:} Hai accesso fisico a spazi che gli
altri non possono raggiungere. Puoi entrare nel Palazzo dei Normanni,
negli uffici dei giustizieri, nei quartieri militari — a patto di
mantenere la copertura.

\textbf{Fardello — Contaminazione:} Hai vissuto tra i francesi troppo
a lungo. Alcuni dei tuoi la pensano come loro, su certi argomenti.
Quando fai qualcosa che metta in pericolo un funzionario angioino
con cui hai rapporti personali, tira Resistere o subisci 1 Stress
dall'esitazione.

\textbf{Legame iniziale:} Il Cospiratore ti ha reclutato tre anni fa.
Non ti sei ancora perdonato di aver accettato.

\filetto

\subsection{Il Mercante}
{\cinzel\footnotesize\color{grigionote} Bartolomeo Lancia}\\
\textit{Il denaro non ha bandiera. Ma tu sì.}

\textbf{Azioni di punta:} Persuadere, Tramare\\
\textbf{Azioni secondarie:} Studiare, Ingannare, Comandare, Manovrare

\textbf{Dono — Rete commerciale:} Hai contatti in ogni porto della
Sicilia e oltre. Puoi richiedere informazioni, merci, o passaggi sicuri
da questi contatti senza tiro — ma ogni richiesta ha un prezzo
narrativo (un favore futuro, una debolezza rivelata).

\textbf{Fardello — Il conto aperto:} Hai finanziato parte della
cospirazione con denaro tuo. Se la rivolta fallisce, sei rovinato.
Questa consapevolezza influenza ogni tua decisione: ogni volta che
un piano sembra vacillare, tira Resistere o agisci in modo egoista.

\textbf{Legame iniziale:} La Spia ti ha fornito documenti falsi
che ti hanno salvato la vita. Non l'hai ancora ripagata.

\filetto

\subsection{Il Soldato Disertore}
{\cinzel\footnotesize\color{grigionote} Drogo de Metz}\\
\textit{Hai combattuto per Carlo d'Angiò. Poi hai visto cosa difendevi.}

\textbf{Azioni di punta:} Combattere, Resistere\\
\textbf{Azioni secondarie:} Agire, Manovrare, Comandare, Sfidare

\textbf{Dono — Disciplina militare:} Quando proteggi un altro PG da
un danno fisico, il danno viene ridotto di 1 casella. Puoi farlo
anche in posizione Disperata senza tiro.

\textbf{Fardello — La faccia conosciuta:} I soldati angioini di Palermo
ti riconoscono. Ogni volta che ti muovi in zone presidiate, il GM può
dichiarare che qualcuno ti ha notato — a meno che tu non abbia preso
precauzioni specifiche.

\textbf{Legame iniziale:} Il Messaggero ti ha convinto a cambiare
parte. Non eri difficile da convincere.

\filetto

\subsection{Il Nobile in Esilio}
{\cinzel\footnotesize\color{grigionote} Tancredi di Chiaramonte}\\
\textit{Discendente di una famiglia che Federico amava. Carlo odia.}

\textbf{Azioni di punta:} Comandare, Persuadere\\
\textbf{Azioni secondarie:} Studiare, Tramare, Resistere, Sfidare

\textbf{Dono — Il nome antico:} Il tuo cognome apre porte tra i baroni
siciliani. Quando ti presenti apertamente, i baroni non possono
ignorarti — anche se non si fidano di te.

\textbf{Fardello — Il bersaglio:} Sei sulla lista di sorveglianza
angioina. I tuoi movimenti vengono notati. Ogni volta che esci
allo scoperto, il GM può aumentare l'allerta angioina di un grado.

\textbf{Legame iniziale:} Il Cospiratore conosce tuo fratello, che
è ancora a Barcellona. Le sue notizie non sono sempre buone.

\filetto

\subsection{Il Falsario}
{\cinzel\footnotesize\color{grigionote} Niccolò il Cretese}\\
\textit{Documenti, sigilli, identità. Tutto ciò che non esiste, tu lo crei.}

\textbf{Azioni di punta:} Ingannare, Agire\\
\textbf{Azioni secondarie:} Studiare, Tramare, Manovrare, Resistere

\textbf{Dono — Il documento giusto:} Una volta per scena, puoi
dichiarare di avere con te un documento che risolve un problema
burocratico o di accesso. Non occorre un tiro — ma il documento
regge solo finché non viene esaminato da un esperto.

\textbf{Fardello — La traccia:} Ogni documento che produci lascia
una firma stilistica. Se le autorità catturano abbastanza tuoi
lavori, possono iniziare a collegarli. Tieni il conto: dopo il
terzo documento usato, il GM può dichiarare che qualcuno ha notato
la somiglianza.

\textbf{Legame iniziale:} Il Mercante è il tuo principale cliente.
Sa troppo di te per essere solo un cliente.

\filetto

\subsection{Il Capitano di Mare}
{\cinzel\footnotesize\color{grigionote} Riccardo della Loggia}\\
\textit{Tre lingue, quattro bandiere, un solo equipaggio. Quando
soffia il vento giusto, sei al servizio di qualcuno.}

\textbf{Azioni di punta:} Manovrare, Comandare\\
\textbf{Azioni secondarie:} Studiare, Persuadere, Ingannare, Resistere

\textbf{Background:} Riccardo capitana una nave latina di mezza
stazza che batte tre rotte: Palermo–Barcellona,
Palermo–Costantinopoli, Palermo–Tunisi. Nato a Trapani, ha
imparato il greco a Negroponte e il catalano a Maiorca. La sua
\textit{Stella di Mar} ha portato lettere cifrate per da Procida,
oro bizantino per i Cospiratori palermitani, e seta legittima
per la corte angioina che lo paga ufficialmente. I tre committenti
non lo sanno l'uno dell'altro. O fanno finta di non saperlo.

\textbf{Dono — La rotta segreta:} Una volta per sessione, puoi
dichiarare di aver stabilito una via veloce e segreta — per
mare, fra due porti — che fa arrivare un carico specifico
(lettere, denaro, una persona) senza intercettazione angioina.
L'arrivo avviene fuori scena; le conseguenze arrivano in scena.

\textbf{Fardello — Il porto sorvegliato:} La flotta di Carlo sta
prendendo forma a Palermo. Ogni volta che usi la rotta segreta,
il GM avanza l'orologio della «sorveglianza portuale». Quando
è pieno, la tua nave viene sequestrata — e non puoi più lasciare
Palermo via mare per il resto della sessione.

\textbf{Legame iniziale:} Il Mercante ti ha pagato tre traversate.
Non ha ancora pagato la quarta. Voi tutti e due sapete perché.

\secsep

%% ============================================================
%% FAZIONE II: BARONI
%% ============================================================
\section{I Baroni Siciliani — Playbook}

\subsection{Il Barone}
{\cinzel\footnotesize\color{grigionote} Simone di Mazzarino}\\
\textit{Hai terra, uomini, e obblighi verso entrambi. E verso tuo figlio.}

\textbf{Azioni di punta:} Comandare, Persuadere\\
\textbf{Azioni secondarie:} Studiare, Tramare, Sfidare, Resistere

\textbf{Dono — La parola data:} Quando stringi un accordo con un altro
personaggio (PG o PNG) e lo dichiari esplicitamente, entrambi
guadagnate +1 dado ai tiri che riguardano quell'accordo, finché
dura. Se uno di voi lo rompe, l'altro guadagna un dado bonus su
qualsiasi azione contro di lui.

\textbf{Fardello — Il peso del feudo:} Le tue decisioni non riguardano
solo te. Ogni volta che scegli un'azione rischiosa, il GM può
chiedere: \textit{e i tuoi uomini? E la tua famiglia?} Devi
rispondere prima di tirare.

\textbf{Legame iniziale:} Il Cavaliere ti è fedele. Speri che questa
fedeltà sopravviva alle scelte che stai per fare.

\filetto

\subsection{Il Cavaliere}
{\cinzel\footnotesize\color{grigionote} Riccardo Bonello}\\
\textit{Hai giurato fedeltà. A chi, esattamente, è diventato complicato.}

\textbf{Azioni di punta:} Combattere, Sfidare\\
\textbf{Azioni secondarie:} Comandare, Resistere, Agire, Manovrare

\textbf{Dono — L'addestramento:} Quando ti trovi in combattimento
e hai almeno un alleato accanto, sei sempre in posizione Rischiosa
(non Disperata) a meno che tu non sia solo e circondato.

\textbf{Fardello — Il giuramento vecchio:} Hai giurato a Carlo
d'Angiò (o a un suo vassallo) in un momento in cui non avevi
scelta. Quel giuramento esiste ancora — giuridicamente, spiritualmente.
Ogni volta che agisci contro gli interessi angioini, tira Resistere
o subisci 1 Stress dalla dissonanza.

\textbf{Legame iniziale:} Il Barone ti protegge. Non sei sicuro
di meritarlo.

\filetto

\subsection{La Castellana}
{\cinzel\footnotesize\color{grigionote} Elisabetta di Paternò}\\
\textit{Tuo marito è morto. Il castello è tuo. Gli uomini di Carlo
pensano di poter trattare con una donna.}

\textbf{Azioni di punta:} Comandare, Tramare\\
\textbf{Azioni secondarie:} Persuadere, Studiare, Resistere, Ingannare

\textbf{Dono — La sottovalutazione:} Una volta per scena, puoi
dichiarare che un avversario ti ha sottovalutato. Il tuo prossimo
tiro contro quell'avversario è in posizione Controllata, anche se
la situazione suggerirebbe altrimenti.

\textbf{Fardello — La visibilità:} Una donna che comanda è visibile.
Le autorità angioine ti sorvegliano con attenzione doppia. Ogni
tiro di Manovrare o Ingannare che faccia in zone controllate dai
francesi parte da posizione Rischiosa, mai Controllata.

\textbf{Legame iniziale:} Il Notaio tiene i tuoi documenti. Sono
documenti che non dovrebbero esistere.

\filetto

\subsection{Il Notaio}
{\cinzel\footnotesize\color{grigionote} Matteo di Termini}\\
\textit{La legge è uno strumento. La domanda è chi la impugna.}

\textbf{Azioni di punta:} Studiare, Tramare\\
\textbf{Azioni secondarie:} Persuadere, Ingannare, Comandare, Resistere

\textbf{Dono — Archivio vivente:} Conosci i precedenti giuridici,
i trattati, le concessioni regie della Sicilia normanna e sveva.
Puoi usare Studiare per trovare una base legale per quasi qualsiasi
azione — il che non la rende sicura, ma la rende giustificabile.

\textbf{Fardello — Il documento scomodo:} Hai tra le mani un documento
che potrebbe cambiare gli equilibri di potere — ma che ti mette
in pericolo solo a possederlo. Il GM sa cosa c'è scritto. Tu non
lo hai ancora letto fino in fondo.

\textbf{Legame iniziale:} La Castellana si fida di te. Hai paura
che questa fiducia sia mal riposta.

\filetto

\subsection{Il Figlio}
{\cinzel\footnotesize\color{grigionote} Goffredo di Mazzarino}\\
\textit{Tuo padre negozia. Tu vorresti agire.}

\textbf{Azioni di punta:} Sfidare, Agire\\
\textbf{Azioni secondarie:} Combattere, Resistere, Aizzare, Manovrare

\textbf{Dono — L'impeto:} Quando agisci per primo in una situazione
di crisi, senza aspettare istruzioni, guadagni +1d. Se la tua azione
funziona, aggiungi +1 al punteggio di fazione.

\textbf{Fardello — L'inesperienza:} La tua posizione di partenza
in qualsiasi situazione che richieda negoziazione o diplomazia
è sempre Rischiosa, mai Controllata. Stai ancora imparando.

\textbf{Legame iniziale:} Il Cavaliere ti ha insegnato a combattere.
Non ti ha ancora insegnato quando non farlo.

\filetto

\subsection{Il Medico}
{\cinzel\footnotesize\color{grigionote} Maestro Filippo Carafa da Salerno}\\
\textit{Curi tutti, anche i francesi. Soprattutto i francesi — sanno cose.}

\textbf{Azioni di punta:} Studiare, Persuadere\\
\textbf{Azioni secondarie:} Resistere, Ingannare, Agire, Invocare

\textbf{Background:} Formato alla Scuola Medica di Salerno, Filippo
Carafa è arrivato a Palermo con una reputazione impeccabile — e qualche
interesse che la Scuola ufficiale non approverebbe. Accanto alla medicina
galenica, ha studiato trattati arabi che circolano solo in ambienti
ristretti: erbe, infusi, preparati che curano dove le erbe normali non
bastano e fanno dormire dove il dolore non lascia pace. Ha trovato in
Yusuf ibn Khalaf al-Balarmì un interlocutore prezioso — sa cose che
nessun salernitano conosce — e in Maimone di Palermo un fornitore
affidabile di materiali che non si trovano nei mercati ufficiali.
Il Barone sa della Scuola. Non sa del resto.

\textbf{Dono — La fiducia del paziente:} Quando hai curato qualcuno,
quella persona ti è debitrice. Puoi richiedere informazioni o un
piccolo favore a chiunque tu abbia curato durante la sessione,
senza tiro — ma una sola volta per persona.

\textbf{Fardello — Il giuramento di Ippocrate:} Non puoi rifiutarti
di curare qualcuno che ne ha bisogno, nemmeno un soldato angioino,
nemmeno un nemico. Se lo fai, subisci 2 Stress immediati.

\textbf{Legami iniziali:} Il Barone ti paga. Ma curi anche i poveri
dell'Albergheria, gratis, la notte. \textit{Legame cross-fazione
(opzionale):} Yusuf ibn Khalaf al-Balarmì (fazione Clero) è il tuo
consulente per la medicina araba — vi scambiate conoscenze da anni.
Maimone di Palermo (fazione Popolo) ti procura ciò che non si trova
nelle spezierie. Entrambi sanno cose di te che il Barone non deve sapere.

\filetto

\subsection{La Moglie del Traditore}
{\cinzel\footnotesize\color{grigionote} Costanza Ventimiglia}\\
\textit{Tuo marito ha ceduto agli Angioini. Tu no.}

\textbf{Azioni di punta:} Resistere, Ingannare\\
\textbf{Azioni secondarie:} Tramare, Persuadere, Studiare, Manovrare

\textbf{Dono — La doppia vita:} Agli occhi delle autorità angioine,
sei la moglie obbediente di un collaboratore. Questo ti dà copertura.
I tuoi tiri di Ingannare nelle zone di controllo angioino
partono sempre da posizione Controllata.

\textbf{Fardello — La casa divisa:} Tuo marito non sa cosa stai
facendo. Se lo scopre, può succedere qualsiasi cosa. Il GM tiene
un orologio segreto: ogni volta che agisci in modo rischioso
vicino a casa, il GM può far avanzare l'orologio. Quando è pieno,
tuo marito scopre tutto.

\textbf{Legame iniziale:} La Castellana è tua cugina. Condividete
più segreti di quanto sia prudente.

\filetto

\subsection{La Matriarca}
{\cinzel\footnotesize\color{grigionote} Donna Margherita Lancia}\\
\textit{Hai sepolto due mariti. Hai sposato sei figlie. Non lasci
la tua casa da undici anni. Eppure tutta Palermo viene a te.}

\textbf{Azioni di punta:} Tramare, Persuadere\\
\textbf{Azioni secondarie:} Studiare, Comandare, Resistere, Invocare

\textbf{Background:} Donna Margherita Lancia ha sepolto suo
marito (caduto a Benevento nel 1266) e il suo primo figlio
(caduto a Tagliacozzo nel 1268). Ha sposato sei figlie in
altrettante case baronali dell'isola. I suoi quattro nipoti
sono ovunque. Da undici anni non lascia il suo \textit{palatium}
alla Vucciria. Non ne ha bisogno: la Sicilia baronale viene a
lei. Riceve nel suo studiolo, due alla volta, senza testimoni.
Sa più cose sulla nobiltà siciliana di quante ce ne siano nei
registri della Magna Curia. Forse perché ne ha scritte molte lei.

\textbf{Dono — Il filo delle case:} Quando parli di una famiglia
baronale con un altro personaggio (PG o PNG), puoi dichiarare
un vincolo specifico — un matrimonio, un debito, un torto — che
li lega alla tua casa. Tu e il GM concordate cosa, ma deve
essere plausibile. Il personaggio deve prenderlo sul serio:
modella come ti parla per il resto della scena.

\textbf{Fardello — La casa attende:} La tua autorità è la tua
casa. Se agisci in modo da metterla pubblicamente in pericolo
(firmando un patto compromettente a tuo nome, prendendo posizione
allo scoperto), il GM avanza l'orologio della successione: i
tuoi figli e nipoti si stanno consultando. Quando è pieno, la
tua casa agisce formalmente contro la tua volontà.

\textbf{Legame iniziale:} Il Notaio custodisce la cassetta dei
tuoi documenti privati. Sa quasi tutto. Il quasi è ciò che conta.

\secsep

%% ============================================================
%% FAZIONE III: CLERO
%% ============================================================
\section{Il Clero — Playbook}

\subsection{Il Vescovo}
{\cinzel\footnotesize\color{grigionote} Berardo di Castronovo}\\
\textit{Servire Dio. Servire Roma. Servire il tuo popolo. Scegli.}

\textbf{Azioni di punta:} Comandare, Invocare\\
\textbf{Azioni secondarie:} Persuadere, Studiare, Tramare, Resistere

\textbf{Dono — L'autorità episcopale:} Una volta per sessione, puoi
emettere un ordine che nessun fedele cattolico (PG o PNG) può
ignorare senza conseguenze spirituali. Il personaggio che disobbedisce
subisce 1 Stress da senso di colpa.

\textbf{Fardello — Roma guarda:} I tuoi movimenti vengono riferiti
alla curia. Se agisci troppo apertamente contro gli interessi di
Martino IV (che è filo-angioino), rischi la rimozione o la scomunica.
Il GM tiene un orologio romano: quando è pieno, arriva una lettera
da Roma.

\textbf{Legame iniziale:} Il Frate Predicatore è il tuo occhi e le
tue orecchie per strada. Non sei sicuro di quanto si fidi di te.

\filetto

\subsection{Il Frate Predicatore}
{\cinzel\footnotesize\color{grigionote} Fra' Tommaso di Caltagirone}\\
\textit{La parola di Dio è un'arma. Dipende solo da chi la impugna.}

\textbf{Azioni di punta:} Aizzare, Invocare\\
\textbf{Azioni secondarie:} Persuadere, Resistere, Studiare, Manovrare

\textbf{Dono — La folla:} Quando parli a un gruppo di persone
(almeno cinque), il tuo tiro di Aizzare ha effetto Standard o
Eccezionale — mai Limitato. Le folle ti ascoltano.

\textbf{Fardello — Il fervore:} La tua capacità di infiammare gli
animi non si spegne a comando. Una volta per sessione, il GM può
dichiarare che hai spinto troppo: la folla ha reagito in modo che
non volevi. Subisci 2 Stress e devi gestire le conseguenze.

\textbf{Legame iniziale:} Il Sacrista ti copre le spalle. Ha visto
cose che non dimentica.

\filetto

\subsection{Il Canonico}
{\cinzel\footnotesize\color{grigionote} Pandolfo di Messina}\\
\textit{Conosci i segreti della cattedrale. Anche quelli che non
sono stati confessati.}

\textbf{Azioni di punta:} Studiare, Tramare\\
\textbf{Azioni secondarie:} Persuadere, Ingannare, Resistere, Invocare

\textbf{Dono — L'archivio:} La cattedrale contiene decenni di
registri, contratti, testamenti, confessioni scritte. Una volta
per sessione, puoi dichiarare di aver trovato un'informazione
negli archivi che è rilevante per la situazione attuale. Il GM
decide cosa contiene, ma deve essere narrativamente utile.

\textbf{Fardello — Il segreto confessionale:} Sai cose che non
puoi rivelare. Alcune di queste cose riguardano persone che sono
in questa storia. Il GM ti dirà, all'inizio della sessione, una
cosa che sai e non puoi dire.

\textbf{Legame iniziale:} Il Vescovo ti ha promosso. Non sei
certo di aver capito perché.

\filetto

\subsection{La Badessa}
{\cinzel\footnotesize\color{grigionote} Suor Adalasia di Troina}\\
\textit{Il tuo convento è un rifugio. Per chi, dipende da te.}

\textbf{Azioni di punta:} Comandare, Resistere\\
\textbf{Azioni secondarie:} Studiare, Persuadere, Tramare, Invocare

\textbf{Dono — Santuario:} Il tuo convento è territorio sacro.
I soldati francesi non possono entrarvi senza violare il diritto
canonico. Puoi nascondere qui persone, documenti, o messaggi —
ma ogni cosa nascosta qui è una responsabilità.

\textbf{Fardello — Le mie sorelle:} Hai decine di donne sotto
la tua protezione. Ogni decisione rischiosa le mette in pericolo.
Prima di ogni azione che potrebbe attirare l'attenzione angioina
sul convento, tira Resistere o subisci 1 Stress dal peso
della responsabilità.

\textbf{Legame iniziale:} Il Frate Predicatore ti porta notizie
dal mondo. Le notizie non sono mai buone abbastanza.

\filetto

\subsection{Il Sacrista}
{\cinzel\footnotesize\color{grigionote} Arnaldo Gennarino}\\
\textit{Controlli le chiavi, le campane, i tempi. Più potere
di quanto sembri.}

\textbf{Azioni di punta:} Agire, Manovrare\\
\textbf{Azioni secondarie:} Studiare, Resistere, Ingannare, Invocare

\textbf{Dono — Le campane:} Puoi suonare le campane di Santo Spirito.
Questo è un atto che tutta la città sente. Usato al momento giusto,
vale automaticamente +1 al punteggio della fazione. Usato al momento
sbagliato, può far scattare una risposta angioina immediata.

\textbf{Fardello — La visibilità del gesto:} Sei l'unico che può
suonare quelle campane. Se qualcosa va storto quella sera, tutti
sapranno che sei stato tu.

\textbf{Legame iniziale:} Il Canonico ti ha insegnato a leggere.
Gli devi più di quello.

\filetto

\subsection{Il Medico Arabo}
{\cinzel\footnotesize\color{grigionote} Yusuf ibn Khalaf al-Balarmì}\\
\textit{La tua famiglia era qui prima dei Normanni. E dopo gli Angioini.}

\textbf{Azioni di punta:} Studiare, Persuadere\\
\textbf{Azioni secondarie:} Agire, Resistere, Ingannare, Invocare

\textbf{Dono — La memoria lunga:} Conosci Palermo più di chiunque
altro — i vicoli, i passaggi, i pozzi. Quando ti muovi nella città,
puoi dichiarare di conoscere una via alternativa. Questo elimina
automaticamente qualsiasi ostacolo fisico al movimento, una volta
per scena.

\textbf{Fardello — Lo straniero in casa propria:} Sei arabo e
cristiano — o arabo e ancora musulmano in privato, il GM decide
con te. In entrambi i casi, sei sospetto agli occhi di tutti.
I tuoi tiri di Comandare con persone che non ti conoscono partono
sempre da posizione Rischiosa.

\textbf{Legame iniziale:} La Badessa ti ha protetto una volta.
Non dimentichi.

\filetto

\subsection{Il Pellegrino}
{\cinzel\footnotesize\color{grigionote} Luca di Bari}\\
\textit{Sei arrivato a Palermo per caso. O forse no.}

\textbf{Azioni di punta:} Manovrare, Ingannare\\
\textbf{Azioni secondarie:} Studiare, Resistere, Persuadere, Invocare

\textbf{Dono — Nessuno mi conosce:} Non sei registrato nelle liste
di sorveglianza angioine. Non sei collegato a nessuna fazione.
I tuoi tiri di Manovrare nelle zone controllate partono sempre
da posizione Controllata.

\textbf{Fardello — Nessuno mi conosce:} Non hai risorse, contatti,
o base sicura. Ogni volta che hai bisogno di un supporto logistico
(un posto dove dormire, un documento, un'arma), devi fare un tiro
per trovarlo.

\textbf{Legame iniziale:} Il Frate Predicatore ti ha dato da mangiare
il primo giorno. Gli hai detto il tuo nome vero. Nessun altro lo sa.

\filetto

\subsection{Il Frate Minore}
{\cinzel\footnotesize\color{grigionote} Fra' Andrea da Caltagirone}\\
\textit{Non hai niente. Non vuoi niente. Per questo, tutti ti
parlano.}

\textbf{Azioni di punta:} Persuadere, Manovrare\\
\textbf{Azioni secondarie:} Resistere, Aizzare, Studiare, Invocare

\textbf{Background:} Fra' Andrea ha preso l'abito francescano sei
anni fa, dopo aver visto suo padre morire di fame ad Augusta.
Cammina senza scarpe in pieno inverno. Mangia ciò che gli si dà.
Non ha residenza fissa: dorme alla porta di chi lo accoglie. La
gente dell'Albergheria gli passa una scodella d'acqua quando lo
vede arrivare. I funzionari francesi non lo arrestano mai — non
perché lo rispettino, ma perché non sanno mai dove sia. Confessa
sulla strada. Battezza i bambini dei quartieri poveri. Non ha mai
chiesto un denaro per niente. Sa cose che il Vescovo non sa.

\textbf{Dono — Povertà come scudo:} Hai meno di chiunque altro.
Non possiedi nulla che si possa rubare o sequestrare. Sei sempre
il benvenuto in qualsiasi casa povera della Sicilia: puoi entrare,
dormire, mangiare, raccogliere informazioni senza tiro nei
quartieri popolari. La gente parla davanti a te perché sei già
parte di loro.

\textbf{Fardello — Il voto:} Non puoi possedere denaro, armi, o
oggetti di valore — anche se ti vengono offerti per la causa.
Se devi maneggiare uno qualsiasi di questi (per portarlo, per
consegnarlo a un altro), subisci 1 Stress per ogni atto. La
purezza della tua vita è un peso reale.

\textbf{Legame iniziale:} Il Frate Predicatore ti rispetta e ti
disprezza in parti uguali. Voi due rappresentate due risposte
diverse alla stessa domanda. Stasera potreste scoprire quale
risposta era giusta.

\secsep

%% ============================================================
%% FAZIONE IV: POPOLO
%% ============================================================
\section{Il Popolo — Playbook}

\subsection{L'Artigiano}
{\cinzel\footnotesize\color{grigionote} Calogero Ferrante}\\
\textit{Fai cose con le mani. Sai anche come rompere cose.}

\textbf{Azioni di punta:} Agire, Combattere\\
\textbf{Azioni secondarie:} Resistere, Manovrare, Sfidare, Aizzare

\textbf{Dono — Gli strumenti del mestiere:} Hai sempre con te
strumenti da lavoro che possono funzionare come armi o come
attrezzi da scasso, senza destare sospetti. Non hai mai un
problema di equipaggiamento fisico.

\textbf{Fardello — Il bottega:} Hai una famiglia che dipende dal
tuo lavoro. Se vieni arrestato o uccisi, loro non mangiano.
Ogni rischio fisico che corri ha un costo emotivo: tira Resistere
o subisci 1 Stress dal pensiero di loro.

\textbf{Legame iniziale:} La Pescivendola è tua vicina. Vi siete
coperti le spalle a vicenda più di una volta.

\filetto

\subsection{La Pescivendola}
{\cinzel\footnotesize\color{grigionote} Perna di Cassaro}\\
\textit{Sei al mercato ogni mattina. Sai tutto di tutti.}

\textbf{Azioni di punta:} Studiare, Persuadere\\
\textbf{Azioni secondarie:} Aizzare, Resistere, Ingannare, Manovrare

\textbf{Dono — Le voci del mercato:} Ogni mattina al mercato,
puoi dichiarare di aver sentito una voce rilevante. Il GM ti
dice una cosa vera sulla situazione attuale — non tutto, ma
qualcosa di utile.

\textbf{Fardello — La lingua:} Non riesci a stare zitta quando
vedi un'ingiustizia. Una volta per sessione, il GM può dichiararti
che hai detto qualcosa ad alta voce che non avresti dovuto. Gestisci
le conseguenze.

\textbf{Legame iniziale:} Il Ragazzo porta le tue ceste. Lo hai
cresciuto come un figlio. Non vuoi che si faccia del male.

\filetto

\subsection{Il Ragazzo}
{\cinzel\footnotesize\color{grigionote} Salvo Russo}\\
\textit{Hai quattordici anni e non hai paura di niente. Dovresti.}

\textbf{Azioni di punta:} Agire, Manovrare\\
\textbf{Azioni secondarie:} Sfidare, Resistere, Aizzare, Studiare

\textbf{Dono — Invisibile:} I soldati francesi non ti guardano —
sei solo un ragazzo. I tuoi tiri di Manovrare in aree presidiate
partono sempre da posizione Controllata.

\textbf{Fardello — Impulsivo:} Quando vedi qualcuno che ami in
pericolo, devi tirare Resistere per non reagire immediatamente.
Se fallisci, agisci — e probabilmente peggiori la situazione.

\textbf{Legame iniziale:} La Pescivendola ti ha tirato su. La ami
come una madre. Questo ti spaventa.

\filetto

\subsection{Il Pescatore}
{\cinzel\footnotesize\color{grigionote} Pietro di Cefalù}\\
\textit{Conosci il mare, il porto, e chi viene e va. Soprattutto chi va.}

\textbf{Azioni di punta:} Manovrare, Studiare\\
\textbf{Azioni secondarie:} Agire, Resistere, Combattere, Ingannare

\textbf{Dono — Il porto:} Hai accesso alla banchina e alle barche.
Puoi far uscire o entrare persone dalla città via mare senza tiro,
a meno che il porto non sia militarizzato. Sai anche quando le navi
di Carlo sono pronte a salpare.

\textbf{Fardello — Il debito:} Devi soldi a qualcuno vicino agli
Angioini. Finché non li ripaghi, sei vulnerabile. Il GM può usare
questo debito come leva narrativa in qualsiasi momento.

\textbf{Legame iniziale:} Il Soldato Disertore (fazione Congiurati)
ha attraversato il porto grazie a te. Non sai ancora se è stato
un errore.

\filetto

\subsection{La Vedova}
{\cinzel\footnotesize\color{grigionote} Giacoma la Nera}\\
\textit{Tuo marito è morto per colpa loro. Non l'hanno nemmeno pagata.}

\textbf{Azioni di punta:} Resistere, Aizzare\\
\textbf{Azioni secondarie:} Persuadere, Studiare, Combattere, Invocare

\textbf{Dono — Il lutto visibile:} Quando parli pubblicamente del
torto subito, il tuo tiro di Aizzare è sempre Eccezionale, mai
Standard o Limitato. Il dolore reale muove le persone.

\textbf{Fardello — I figli:} Hai tre figli piccoli. Sono la tua
forza e la tua debolezza. Il GM può mettere uno di loro in pericolo
in qualsiasi momento per creare pressione narrativa.

\textbf{Legame iniziale:} L'Artigiano era amico di tuo marito.
Ti ha promesso di tenerti d'occhio. Non sai se riesce a mantenere
la promessa.

\filetto

\subsection{Il Guaritore di Strada}
{\cinzel\footnotesize\color{grigionote} Maimone di Palermo}\\
\textit{Non sei un medico. Ma sai come tenere in vita la gente.}

\textbf{Azioni di punta:} Agire, Studiare\\
\textbf{Azioni secondarie:} Persuadere, Resistere, Manovrare, Invocare

\textbf{Dono — Mani sicure:} Quando curi qualcuno (riduci le
caselle di danno), il tuo tiro è sempre almeno Standard. Non
puoi peggiorare la situazione fisica di un paziente con un fallimento.

\textbf{Fardello — La visione:} Vedi troppo. Ogni volta che
sei vicino a qualcuno ferito, subisci 1 Stress dal peso di
ciò che vedi. Questo è il costo di cure tante persone.

\textbf{Legame iniziale:} La Vedova è la tua vicina di casa.
I suoi figli ti chiamano zio.

\filetto

\subsection{Il Vecchio}
{\cinzel\footnotesize\color{grigionote} Basilio di Monreale}\\
\textit{Hai vissuto abbastanza da sapere come va a finire.
E vai avanti lo stesso.}

\textbf{Azioni di punta:} Persuadere, Resistere\\
\textbf{Azioni secondarie:} Studiare, Aizzare, Invocare, Comandare

\textbf{Dono — La memoria:} Ricordi la Sicilia prima degli Angioini.
E prima dei Normanni. Quando dici a qualcuno com'era, puoi usare
Persuadere per ispirare azioni che non avresti pensato possibili.
Una volta per sessione, puoi dare a un altro PG +1d al suo prossimo tiro.

\textbf{Fardello — Il corpo:} Sei vecchio. I tiri di Agire e
Combattere partono sempre da posizione Rischiosa, mai Controllata.
Il tuo corpo non segue più quello che vuoi fare.

\textbf{Legame iniziale:} Il Ragazzo ti ricorda qualcuno. Non gli
dici chi. Forse è per questo che gli vuoi bene.

\filetto

\subsection{L'Ostessa}
{\cinzel\footnotesize\color{grigionote} Sabella d'Albergheria}\\
\textit{Servi vino acido a tutti — francesi, congiurati, soldati,
predicatori. Ascolti tutto. Non dici niente. La porta sul retro
non la chiudi mai.}

\textbf{Azioni di punta:} Studiare, Persuadere\\
\textbf{Azioni secondarie:} Ingannare, Resistere, Manovrare, Aizzare

\textbf{Background:} La taverna di Sabella sta nel vicolo che
sale verso la Cattedrale, ai margini dell'Albergheria. Vino
acidulo del Monte San Giuliano, pane raffermo riscaldato sulla
brace, due tavoli e una panca lunga. Aperta dal vespro all'alba.
Quando i soldati francesi vengono, si siedono al tavolo vicino
alla porta — Sabella li serve in silenzio. Quando gli artigiani
del Cassaro vengono per riunioni che non vogliono in nessun
verbale, si siedono nel retro. Lei chiude la porta principale.
Non chiude la porta sul retro. Sa chi è venuto, quando, con chi,
e quanto ha bevuto. Non dimentica.

\textbf{Dono — La porta sul retro:} La tua taverna è territorio
neutrale. Una volta per scena, puoi dichiarare che due personaggi
(PG o PNG) che non si parlerebbero altrove si sono ritrovati al
tuo tavolo di fondo. La conversazione avviene in posizione
narrativamente neutra — né minaccia né alleanza — finché uno
dei due non la rompe.

\textbf{Fardello — Servi a tutti, alleata di nessuno:} Per restare
aperta, devi servire chiunque paghi. Se prendi posizione apertamente
per una parte, il GM (su richiesta dell'altra) può chiudere la
tua taverna. Se chiude, perdi il tuo strumento principale per
il resto della sessione — e con la taverna perde anche il Popolo
uno dei suoi pochi spazi.

\textbf{Legame iniziale:} L'Artigiano viene ogni giovedì sera.
Ti deve tre mesi di vino. In cambio ti ha sistemato la serratura
del retro. Tu sai che quella serratura si apre con un grimaldello
specifico. Lui non sa che tu lo sai.

\filettodoppio
